% Mechanistic Validity — NEMI 2026 poster, Option C
% Practical/tool framing: "Here's how to audit your own claims."
% Text and tables only, no figures.
%
%   cd poster && lualatex mechval_poster_option_c.tex

\documentclass[28pt, a0paper, portrait]{tikzposter}

\usepackage[T1]{fontenc}
\usepackage{amsmath,amssymb}
\usepackage{array}
\usepackage{booktabs}
\usepackage{enumitem}
\usepackage{microtype}

\tikzposterlatexaffectionproofoff

% ── Color palette ───────────────────────────────────────
\definecolor{darktext}{HTML}{2C3E50}
\definecolor{blockbg}{HTML}{FFFFFF}
\definecolor{tierTriangulated}{HTML}{27AE60}
\definecolor{tierSupported}{HTML}{2980B9}
\definecolor{tierCausal}{HTML}{F39C12}
\definecolor{tierProposed}{HTML}{95A5A6}
\definecolor{tierDisconfirmed}{HTML}{E74C3C}
\definecolor{construct}{HTML}{C0392B}
\definecolor{newlabel}{HTML}{8E44AD}
\definecolor{practicalgreen}{HTML}{27AE60}

% ── Theme ───────────────────────────────────────────────
\usetheme{Default}

\definecolorstyle{MechValColors}{
  \definecolor{colorOne}{HTML}{2C3E50}
  \definecolor{colorTwo}{HTML}{3498DB}
  \definecolor{colorThree}{HTML}{ECF0F1}
}{
  \colorlet{backgroundcolor}{colorThree}
  \colorlet{framecolor}{colorOne}
  \colorlet{titlefgcolor}{white}
  \colorlet{titlebgcolor}{colorOne}
  \colorlet{blocktitlebgcolor}{colorTwo}
  \colorlet{blocktitlefgcolor}{white}
  \colorlet{blockbodybgcolor}{blockbg}
  \colorlet{blockbodyfgcolor}{darktext}
  \colorlet{notefgcolor}{darktext}
  \colorlet{notebgcolor}{tierCausal!15}
  \colorlet{notefrcolor}{tierCausal}
}
\usecolorstyle{MechValColors}

\defineblockstyle{MechValBlock}{
  titlewidthscale=1.0, bodywidthscale=1.0, titleleft,
  titleoffsetx=0pt, titleoffsety=0pt,
  bodyoffsetx=0pt, bodyoffsety=0pt, bodyverticalshift=0pt,
  roundedcorners=12, linewidth=2pt, titleinnersep=10pt, bodyinnersep=12pt
}{
  \draw[rounded corners=\blockroundedcorners, color=blocktitlebgcolor,
        fill=blockbodybgcolor, line width=\blocklinewidth]
    (blockbody.south west) rectangle (blocktitle.north east);
  \ifBlockHasTitle
    \draw[rounded corners=\blockroundedcorners,
          fill=blocktitlebgcolor, color=blocktitlebgcolor,
          line width=\blocklinewidth]
      (blocktitle.south west) rectangle (blocktitle.north east);
  \fi
}
\useblockstyle{MechValBlock}

\definetitlestyle{MechValTitle}{
  width=\paperwidth, roundedcorners=0, linewidth=0pt, innersep=20pt,
  titletotopverticalspace=12mm, titletoblockverticalspace=20mm
}{
  \draw[fill=titlebgcolor, color=titlebgcolor]
    (\titleposleft, \titleposbottom)
    rectangle (\titleposright, \titlepostop);
}
\usetitlestyle{MechValTitle}

\title{%
  \parbox{0.92\linewidth}{\centering\bfseries
    Mechanistic Validity\\[6pt]
    A Validity Audit for Mechanistic Interpretability}}
\author{Elliot Tower\quad\texttt{elliot@elliottower.ai}}
\institute{}

\begin{document}

\maketitle[width=\paperwidth]

% ════════════════════════════════════════════════════════
% ROW 1 — GAP + FIVE TYPES + EVIDENCE CARD
% ════════════════════════════════════════════════════════
\begin{columns}
\column{0.33}

\block{The Gap}{
  Every field that makes mechanistic claims has validity
  standards.

  \vspace{4mm}
  \begin{itemize}[leftmargin=*, itemsep=3mm]
    \item \textbf{Neuroscience}: Craver's constitutive
      relevance --- mutual manipulability at a declared level
    \item \textbf{Genetics}: complementation tests, rescue
      experiments, sensitivity analysis
    \item \textbf{Pharmacology}: dose-response curves,
      the ORBITA trial
    \item \textbf{Measurement theory}: Cronbach and Meehl's
      construct validity (1955)
  \end{itemize}

  \vspace{4mm}
  Mechanistic interpretability borrows the interventional
  vocabulary --- ablation, patching, knockout --- but not the
  statistical discipline that comes with it.

  \vspace{4mm}
  We adapt \textbf{five validity types} for neural network
  mechanisms. \textbf{36 criteria}. Applied to
  \textbf{16 published claims}.
}

\column{0.34}

{
\colorlet{blocktitlebgcolor}{construct}
\block{Five Types, One Dependency Chain}{
  Each type answers one question. The chain means
  \textbf{construct is a ceiling on all the rest}.

  \vspace{5mm}
  \renewcommand{\arraystretch}{1.5}
  {\large
  \begin{tabular}{@{}r>{\raggedright\arraybackslash}p{0.75\linewidth}@{}}
    \textbf{Construct} & Is the circuit a natural kind?
      {\small (C4: discriminant validity)} \\
    $\downarrow$ & \\[-6mm]
    \textbf{Measurement} & Does the metric measure what
      it claims? {\small (M2: baseline separation)} \\
    $\downarrow$ & \\[-6mm]
    \textbf{Internal} & Is the causal claim warranted?
      {\small (I4: specificity)} \\
    $\downarrow$ & \\[-6mm]
    \textbf{External} & Does the result generalize?
      {\small (E1: intervention reach)} \\
    $\downarrow$ & \\[-6mm]
    \textcolor{newlabel}{\textbf{Interpretive}} &
      \textcolor{newlabel}{Is the label justified?}
      {\small \textcolor{newlabel}{(V2: level-evidence match)
      \textbf{[NEW]}}} \\
  \end{tabular}}

  \vspace{5mm}
  {\Large The verdict is the \textbf{minimum}, not the average.}

  \vspace{3mm}
  No amount of causal evidence rescues a claim whose target
  concept is incoherent. No measurement reliability
  compensates for measuring the wrong thing.
}
}

\column{0.33}

\block{How to Read a Verdict}{
  Every claim gets an \textbf{evidence card}: the atomic
  reporting unit linking measurement to verdict.

  \vspace{4mm}
  \renewcommand{\arraystretch}{1.35}
  \begin{tabular}{@{}p{0.25\linewidth}p{0.68\linewidth}@{}}
    \toprule
    \textbf{Claim} & A 26-head circuit implements IOI
      in GPT-2 Small \\[2pt]
    \textbf{Claimed} & Triangulated \\[2pt]
    \textbf{Supported} & \textcolor{tierCausal}{\textbf{Causally
      Suggestive}} \\[2pt]
    \textbf{Capped by} & Method-conditional faithfulness:
      87\% under mean ablation, below 0\% under others \\[2pt]
    \textbf{To close} & Specificity test, double
      dissociation, cross-model replication \\
    \bottomrule
  \end{tabular}

  \vspace{5mm}
  Three questions for any mechanistic claim:
  \begin{enumerate}[leftmargin=*, itemsep=2mm]
    \item What \textbf{description mode} is declared?
    \item What \textbf{tier} does the evidence support?
    \item What \textbf{caps} it?
  \end{enumerate}
}

\end{columns}

% ════════════════════════════════════════════════════════
% ROW 2 — 16-CLAIM TABLE + DIAGNOSTIC PATTERNS
% ════════════════════════════════════════════════════════
\begin{columns}
\column{0.67}

{
\colorlet{blocktitlebgcolor}{colorOne}
\block{Sixteen Claims from Fifteen Papers}{
  {\small
  \renewcommand{\arraystretch}{1.25}
  \begin{tabular}{@{}>{\raggedright\arraybackslash}p{3.6cm}
                    >{\raggedright\arraybackslash}p{7.0cm}
                    >{\raggedright\arraybackslash}p{5.5cm}@{}}
    \toprule
    \textbf{Verdict} & \textbf{Claim} & \textbf{Scope} \\
    \midrule
    \textcolor{tierTriangulated}{\textbf{Triangulated}}
      & Induction heads: token copying
      & 1L--70B, copying only \\
    \midrule
    \textcolor{tierSupported}{\textbf{Mech.\ Supported}}
      & Fourier mechanism (modular addition)
      & 1L ReLU transformer \\
      & Greater-Than circuit
      & GPT-2 Small \\
      & Copy Suppression (L10H7)
      & GPT-2 Small \\
      & Superposition (toy networks)
      & 4 ReLU nets \\
      & Steering vector (refusal)
      & 13 models, 5 families \\
      & Global Workspace / J-space
      & 4 Claude models \\
      & Successor Heads
      & Pythia-1.4B \\
    \midrule
    \textcolor{tierCausal}{\textbf{Causally Suggestive}}
      & Othello World Model
      & Othello-GPT \\
      & Docstring Circuit
      & 4L attn-only \\
      & Gender Bias Circuits
      & GPT-2 + 5 families \\
      & IOI Circuit
      & GPT-2 Small \\
    \midrule
    \textcolor{tierProposed}{\textbf{Proposed}}
      & Probing Classifiers
      & Review (no single task) \\
    \midrule
    \textcolor{tierDisconfirmed}{\textbf{Disconfirmed}}
      & Knowledge Neurons
      & BERT-base \\
      & SAE features as canonical units
      & Pythia \\
      & Induction heads: general ICL
      & $>$1B models \\
    \bottomrule
  \end{tabular}}

  \vspace{4mm}
  No claim reaches \textbf{Validated}. The highest tier is
  Triangulated --- by the field's oldest mechanism, and only
  the narrow reading.
}
}

\column{0.33}

{
\colorlet{blocktitlebgcolor}{tierDisconfirmed}
\block{Three Diagnostic Patterns}{
  These recur across claims:

  \vspace{4mm}
  \textbf{1. Construct ceiling.}
  The Othello ``world model'' label asserts more than
  ``linearly decodable board state.'' The intervention
  supports the weaker reading; the label claims the stronger
  one. V2 caps the verdict.

  \vspace{4mm}
  \textbf{2. Method conditioning.}
  IOI faithfulness swings from below 0\% to over 100\%
  across six methodological choices. The causal claim is
  method-conditional in both directions. E1 caps the verdict.

  \vspace{4mm}
  \textbf{3. Same paper, two verdicts.}
  Induction heads for token copying: \textbf{Triangulated}.
  Induction heads for general in-context learning:
  \textbf{Disconfirmed}. The authors anticipated this ---
  their Argument~6 concedes the case ``is stronger for small
  models.'' One paper, two claims, two tiers.
}
}

\end{columns}

% ════════════════════════════════════════════════════════
% ROW 3 — PRACTICAL + LIMITATIONS + TAKEAWAY
% ════════════════════════════════════════════════════════
\begin{columns}
\column{0.33}

{
\colorlet{blocktitlebgcolor}{practicalgreen}
\block{For Your Next Paper}{
  \begin{enumerate}[leftmargin=*, itemsep=4mm]
    \item \textbf{Declare the description mode.}
      Implementational? Functional? Topographic? The mode
      determines what evidence the claim owes.

    \item \textbf{Test specificity, not just necessity.}
      What does your circuit \emph{not} do? Ablation shows
      the component matters; specificity shows it matters
      \emph{for this}.

    \item \textbf{Run a baseline.} Does a random subspace
      of matched rank achieve comparable scores? A metric
      without a baseline is a metric without meaning.

    \item \textbf{Report what caps your tier.} Name the
      criterion, not just the faithfulness number. The gap
      between ``claimed'' and ``supported'' is the distance
      to close.
  \end{enumerate}
}
}

\column{0.34}

\block{Limitations}{
  \begin{itemize}[leftmargin=*, itemsep=3mm]
    \item The framework has \textbf{not been validated
      against itself}. Three planned checks are specified
      and pre-registered but not run: ground-truth
      calibration, prospective audit, inter-rater agreement.

    \item \textbf{Expect poor agreement.} Annotators
      rating evidence--claim gaps across 186 claims agree
      at $\alpha = 0.11$, near chance. Autointerpretability
      scores show intraclass correlations as low as $-0.049$.

    \item \textbf{Scoring is retrospective}, by the
      framework's author. This is the arrangement the
      prospective audit exists to fix.

    \item \textbf{Verdicts are claim-level, not
      paper-level.} Several papers scored below Triangulated
      state their own scope correctly. The cap reflects what
      evidence licenses, not care taken.
  \end{itemize}
}

\column{0.33}

{
\colorlet{blocktitlebgcolor}{colorOne}
\block{Takeaway}{
  {\Large\bfseries Report the tier your evidence supports,
  and name the criterion that caps it.}

  \vspace{5mm}
  \begin{itemize}[leftmargin=*, itemsep=3mm]
    \item \textbf{36} criteria across \textbf{5} validity
      types in dependency order
    \item \textbf{16} claims audited from \textbf{15} papers
    \item \textbf{0} claims reach Validated
    \item \textbf{1} reaches Triangulated --- the narrow
      reading of induction heads
    \item \textbf{3} Disconfirmed, on evidence the original
      authors anticipated
  \end{itemize}

  \vspace{6mm}
  \begin{center}
  \texttt{elliot@elliottower.ai}

  \vspace{3mm}
  {\small Framework, criteria, and all sixteen evidence cards:}

  \texttt{mechanisticresearch.org}
  \end{center}
}
}

\end{columns}

\end{document}
